% Λύση του 70ού προβλήματος της ενότητας «Άλυτα Προβλήματα».
\subsection*{Λύση Προβλήματος 70 --- Υποχώρηση παροχής μετά από βροχή}

\begin{alist}
  \item Έχουμε
        \[
          Q(t)=80e^{-t/6}.
        \]
        Παραγωγίζουμε:
        \[
          Q'(t)=80\cdot\left(-\frac{1}{6}\right)e^{-t/6}
          =
          -\frac{40}{3}e^{-t/6}.
        \]
        Επειδή
        \[
          e^{-t/6}>0
        \]
        για κάθε \(t\), έχουμε
        \[
          Q'(t)<0.
        \]
        Άρα η παροχή είναι γνησίως φθίνουσα στο διάστημα \([0,24]\).

  \item Θέλουμε
        \[
          Q(t)=40.
        \]
        Δηλαδή
        \[
          80e^{-t/6}=40.
        \]
        Διαιρούμε με \(80\):
        \[
          e^{-t/6}=\frac{1}{2}.
        \]
        Παίρνουμε λογάριθμο:
        \[
          -\frac{t}{6}=\ln\frac{1}{2}.
        \]
        Επειδή \(\ln\frac{1}{2}=-\ln2\), προκύπτει
        \[
          t=6\ln2.
        \]
        Άρα
        \[
          {t=6\ln2\ \text{ώρες}}
        \]
        δηλαδή περίπου
        \[
          {t\approx4{,}16\ \text{ώρες}}.
        \]

  \item Ο ζητούμενος όγκος είναι
        \[
          V=3600\int_0^{12}80e^{-t/6}\,dt.
        \]
        Υπολογίζουμε πρώτα το ολοκλήρωμα:
        \[
          \int_0^{12}80e^{-t/6}\,dt
          =
          \left[-480e^{-t/6}\right]_0^{12}.
        \]
        Άρα
        \[
          \int_0^{12}80e^{-t/6}\,dt
          =
          -480e^{-2}+480
          =
          480(1-e^{-2}).
        \]
        Επομένως
        \[
          V=3600\cdot480(1-e^{-2}).
        \]
        Δηλαδή
        \[
          {V=1\,728\,000(1-e^{-2})\si{m^3}}.
        \]
        Αριθμητικά,
        \[
          {V\approx1\,494\,000\si{m^3}}.
        \]
\end{alist}

\paragraph{Πηγές και έλεγχος ρεαλιστικότητας}
Το μοντέλο γραμμικής δεξαμενής χρησιμοποιείται στην υδρολογία για την
περιγραφή απορροής. Όταν δεν υπάρχει νέα βροχή ή άλλη τροφοδότηση, η
παροχή μπορεί να περιγραφεί με εκθετική μείωση
\[
  Q(t)=Q_0e^{-at}
\]
\cite{linearReservoirRunoffModel}.

Στο πρόβλημα πήραμε \(Q_0=80\si{m^3/s}\) και χρόνο υποχώρησης
\(6\) ώρες. Η τιμή \(Q(0)=80\si{m^3/s}\) μπορεί να αντιστοιχεί σε
αυξημένη παροχή μικρού ή μεσαίου ποταμού αμέσως μετά από έντονη
βροχόπτωση, ενώ
\[
  Q(24)=80e^{-4}\approx1{,}47\si{m^3/s}
\]
δείχνει ότι η πλημμυρική απορροή έχει σχεδόν εκτονωθεί μετά από μία
ημέρα. Ο παράγοντας \(3600\) στον τύπο του όγκου μπαίνει επειδή το
\(t\) μετριέται σε ώρες, ενώ η παροχή δίνεται σε
\(\si{m^3/s}\).
